% Part: lambda-calculus
% Chapter: lambda-definability
% Section: arithmetical-functions

\documentclass[../../../include/open-logic-section]{subfiles}

\begin{document}

\olfileid{lam}{rep}{arf}
\olsection{الدوال الحسابية \usetoken{S}{lambda definable}}

\begin{prop}
  \ollabel{prop:succ-ld}
  دالة الخلف~$\Succ$ من الدوال !!{lambda definable}.
\end{prop}

\begin{proof}
نختار الحد الذي !!{lambda define}s دالة الخلف كما يلي:
\begin{align*}
  \fn{Succ} & \ident \lambd[\OLId{latin-ordinary}{a}][\lambd[fx][\OLId{latin-ordinary}{f} ({\OLId{latin-ordinary}{a}} \OLId{latin-ordinary}{f} \OLId{latin-ordinary}{x})]].
  \intertext{ومقتضى اصطلاحات الاختصار أن هذا هو}
  \fn{Succ} & \ident \lambd[\OLId{latin-ordinary}{a}][\lambd[\OLId{latin-ordinary}{f}][\lambd[\OLId{latin-ordinary}{x}][(\OLId{latin-ordinary}{f} (({\OLId{latin-ordinary}{a}} \OLId{latin-ordinary}{f}) \OLId{latin-ordinary}{x}))]]].
  \intertext{تأخذ~$\fn{Succ}$ عددًا~$a$، فتعطي دالة أخرى هي~$\lambd[fx][f ({a} f x)]$.
    وليست هذه الدالة في ذات كتابتها عدد تشيرش. لكنها تختزل إلى عدد تشيرش
    متى كان~$a$ عدد تشيرش. وبيانه:}
   (\lambd[\OLId{latin-ordinary}{a}][\lambd[fx][\OLId{latin-ordinary}{f} ({\OLId{latin-ordinary}{a}} \OLId{latin-ordinary}{f} \OLId{latin-ordinary}{x})]])\,\num{\OLId{latin-ordinary}{n}} & \redone
   \lambd[fx][\OLId{latin-ordinary}{f} ({\num{\OLId{latin-ordinary}{n}}} \OLId{latin-ordinary}{f} \OLId{latin-ordinary}{x})].
   \intertext{وفي الداخل يقع~$\num{n}fx$ موضع اختزال؛ فإن~$\num{n}$
     هو $\lambd[fx][f^nx]$. ومن ثم~$\num{n}fx \red f^nx$،
     فيكون اختزال الحد كله}
   \fn{Succ}\, \num \OLId{latin-ordinary}{n} & \red \lambd[fx][\OLId{latin-ordinary}{f}(\OLId{latin-ordinary}{f}^\OLId{latin-ordinary}{n}(\OLId{latin-ordinary}{x}))],
\end{align*}
وهذه هي~$\num{\OLId{latin-ordinary}{n}+\OLMathNumeral{1}}$.
\end{proof}

\begin{ex}
  ولنبسط ما يجري عند تطبيق~$\fn{Succ}$ في~$\num{\OLMathNumeral{0}}$،
  أي في~$\lambd[fx][\OLId{latin-ordinary}{x}]$، بإثبات الحدود تامة:
  \begin{align*}
    \fn{Succ}\, \num{\OLMathNumeral{0}} & \ident (\lambd[\OLId{latin-ordinary}{a}][\lambd[\OLId{latin-ordinary}{f}][\lambd[\OLId{latin-ordinary}{x}][(\OLId{latin-ordinary}{f} (({\OLId{latin-ordinary}{a}} \OLId{latin-ordinary}{f}) \OLId{latin-ordinary}{x}))]]])(\lambd[\OLId{latin-ordinary}{f}][\lambd[\OLId{latin-ordinary}{x}][\OLId{latin-ordinary}{x}]])\\
    & \redone \lambd[\OLId{latin-ordinary}{f}][\lambd[\OLId{latin-ordinary}{x}][(\OLId{latin-ordinary}{f} (({(\lambd[\OLId{latin-ordinary}{f}][\lambd[\OLId{latin-ordinary}{x}][\OLId{latin-ordinary}{x}]])} \OLId{latin-ordinary}{f}) \OLId{latin-ordinary}{x}))]] \\
    & \redone \lambd[\OLId{latin-ordinary}{f}][\lambd[\OLId{latin-ordinary}{x}][(\OLId{latin-ordinary}{f} ({(\lambd[\OLId{latin-ordinary}{x}][\OLId{latin-ordinary}{x}])} \OLId{latin-ordinary}{x}))]] \\
    & \redone \lambd[\OLId{latin-ordinary}{f}][\lambd[\OLId{latin-ordinary}{x}][(\OLId{latin-ordinary}{f} \OLId{latin-ordinary}{x})]] \ident \num{\OLMathNumeral{1}}\\
  \end{align*}
\end{ex}

\begin{prob}
  هذا الحد
  \begin{align*}
    \fn{Succ}' & \ident \lambd[{\OLId{latin-ordinary}{n}}][\lambd[fx][{\OLId{latin-ordinary}{n}} \OLId{latin-ordinary}{f} (\OLId{latin-ordinary}{f} \OLId{latin-ordinary}{x})]]
  \end{align*}
  !!{lambda define}s دالة الخلف أيضًا؛ فبيّن وجه ذلك.
\end{prob}

\begin{prop}
  \ollabel{prop:add-ld}
  دالة الجمع~$\Add$ من الدوال !!{lambda definable}.
\end{prop}

\begin{proof}
لنا أن !!{lambda define} الجمع بأحد الحدين:
\begin{align*}
  \fn{Add} & \ident \lambd[{\OLId{latin-ordinary}{a}}{\OLId{latin-ordinary}{b}}][\lambd[fx][{\OLId{latin-ordinary}{a}} \OLId{latin-ordinary}{f} ({\OLId{latin-ordinary}{b}} \OLId{latin-ordinary}{f} \OLId{latin-ordinary}{x})]]
\intertext{أو بالحد الآخر}
  \fn{Add}' & \ident \lambd[{\OLId{latin-ordinary}{a}}{\OLId{latin-ordinary}{b}}][{\OLId{latin-ordinary}{a}}\, \fn{Succ} \, {\OLId{latin-ordinary}{b}}].
\intertext{أما الأول فبيانه: تأخذ~$\fn{Add}$ العددين~$ {a}$ و$ {b}$.
  ثم تعطي دالة تأخذ~$f$ و$x$ وتردّ~$af(bfx)$.
  فإذا كان~$a$ و$b$ عددي تشيرش~$\num{n}$ و$\num{m}$،
  اختزل الناتج إلى $f^{n+m}(x)$، أي إلى $f^{n}(f^{m}(x))$ نفسه.
  وتفصيل الاختزال:}
  (\lambd[{\OLId{latin-ordinary}{a}}{\OLId{latin-ordinary}{b}}][\lambd[fx][{\OLId{latin-ordinary}{a}} \OLId{latin-ordinary}{f} ({\OLId{latin-ordinary}{b}} \OLId{latin-ordinary}{f} \OLId{latin-ordinary}{x})]])\num \OLId{latin-ordinary}{n}\,\num \OLId{latin-ordinary}{m} & \red
  \lambd[fx][\num{\OLId{latin-ordinary}{n}}\, \OLId{latin-ordinary}{f} (\num {\OLId{latin-ordinary}{m}}\, \OLId{latin-ordinary}{f} \OLId{latin-ordinary}{x})] \\
  & \red \lambd[fx][\num{\OLId{latin-ordinary}{n}}\, \OLId{latin-ordinary}{f} (\OLId{latin-ordinary}{f}^\OLId{latin-ordinary}{m} \OLId{latin-ordinary}{x})] \\
  & \red \lambd[fx][\OLId{latin-ordinary}{f}^\OLId{latin-ordinary}{n} (\OLId{latin-ordinary}{f}^\OLId{latin-ordinary}{m} \OLId{latin-ordinary}{x})] \ident \num {\OLId{latin-ordinary}{n}+\OLId{latin-ordinary}{m}}.
\intertext{وأما التمثيل الثاني للجمع~$\fn{Add'}$ فطريقه غير ذلك.
  إذ نطبقه في عددي تشيرش~$\num{n}$ و$\num{m}$ فنحصل على}
\fn{Add}' \num \OLId{latin-ordinary}{n} \,\num \OLId{latin-ordinary}{m}
& \red \num{\OLId{latin-ordinary}{n}}\, \fn{Succ}\, \num{\OLId{latin-ordinary}{m}}.
\intertext{والأصل أن~$\num{n} f x$ يختزل دائمًا إلى~$f^n(x)$. فيلزم}
  \num{\OLId{latin-ordinary}{n}}\,  \fn{Succ}\, \num{\OLId{latin-ordinary}{m}} & \red \fn{Succ}^\OLId{latin-ordinary}{n}(\num{\OLId{latin-ordinary}{m}}).
\end{align*}
ثم إن~$\fn{Succ}$ !!{lambda define}s دالة الخلف؛ وتكرار دالة
الخلف~$\OLId{latin-ordinary}{n}$ مرة ابتداءً من~$\OLId{latin-ordinary}{m}$ يعطي~$\OLId{latin-ordinary}{n}+\OLId{latin-ordinary}{m}$.
فآخر الاختزال إذن هو~$\num{\OLId{latin-ordinary}{n}+\OLId{latin-ordinary}{m}}$.
\end{proof}

\begin{prop}
  \ollabel{prop:mult-ld}
  الضرب !!{lambda definable} بهذا الحد:
  \[
  \fn{Mult} \ident \lambd[ab][\lambd[fx][\OLId{latin-ordinary}{a} (\OLId{latin-ordinary}{b} \OLId{latin-ordinary}{f}) \OLId{latin-ordinary}{x}]]
  \]
\end{prop}

\begin{proof}
  لنتتبع تطبيق~$\fn{Mult}$ في عددي تشيرش~$\num{\OLId{latin-ordinary}{n}}$ و$\num{\OLId{latin-ordinary}{m}}$.
  يختزل~$\fn{Mult} \, \num{\OLId{latin-ordinary}{n}} \, \num{\OLId{latin-ordinary}{m}}$
  إلى~$\lambd[fx][\num{\OLId{latin-ordinary}{n}}(\num{\OLId{latin-ordinary}{m}}\, \OLId{latin-ordinary}{f})\OLId{latin-ordinary}{x}]$.
  والحد~$\num{\OLId{latin-ordinary}{m}} \OLId{latin-ordinary}{f}$ يعرّف دالة تكرر تطبيق~$\OLId{latin-ordinary}{f}$ في وسيطها~$\OLId{latin-ordinary}{m}$ مرة.
  ولذلك يكون فعل~$\num{\OLId{latin-ordinary}{n}} (\num{\OLId{latin-ordinary}{m}} \OLId{latin-ordinary}{f}) \OLId{latin-ordinary}{x}$
  أن يكرر الدالة «طبّق~$\OLId{latin-ordinary}{f}$ بمقدار~$\OLId{latin-ordinary}{m}$ من المرات»
  نفسها~$\OLId{latin-ordinary}{n}$ مرة، بادئًا من~$\OLId{latin-ordinary}{x}$.
  فكأننا طبقنا~$\OLId{latin-ordinary}{f}$ في~$\OLId{latin-ordinary}{x}$ بمقدار~$\OLId{latin-ordinary}{n}\cdot \OLId{latin-ordinary}{m}$ من المرات؛
  والصورة السوية لذلك هي عدد تشيرش~$\num{nm}$ بعينه.
\end{proof}

\begin{editorial}
ويجوز مزيد من التبسيط باختزال~$\eta$ فنكتب
\[
  \fn{Mult} \ident \lambd[ab][\lambd[\OLId{latin-ordinary}{f}][\OLId{latin-ordinary}{a} (\OLId{latin-ordinary}{b} \OLId{latin-ordinary}{f})]].
\]

غير أن هذا يقتضي تقديم شرح اختزال~$\eta$ أولًا.
\end{editorial}

\begin{prob}
لنا أن !!{lambda define} الضرب بالحد
\[
  \fn{Mult}' \ident \lambd[ab][\OLId{latin-ordinary}{a} (\fn{Add}\, \OLId{latin-ordinary}{b}) \num{\OLMathNumeral{0}}].
\]
بيّن كيف يؤدي هذا الحد المطلوب.
\end{prob}

وأما رفع العدد إلى قوة بحد~$\lambd$، إذا كان الأس موجبًا، فيمثله هذا الحد الموجز:
\[
  \fn{Exp} \ident \lambd[be][\OLId{latin-ordinary}{e} \OLId{latin-ordinary}{b}].
\]
فيه الوسيط الأول~$\OLId{latin-ordinary}{b}$ هو الأساس، والثاني~$\OLId{latin-ordinary}{e}$ هو الأس.
ووجهه الحدسي أن~$\OLId{latin-ordinary}{e} \OLId{latin-ordinary}{f}$ هو~$\OLId{latin-ordinary}{f}^\OLId{latin-ordinary}{e}$ بمقتضى ترميز الأعداد الذي اخترناه.
فإن خفي هذا الوجه، فطريق آخر لرفع العدد إلى قوة هو الضرب المتكرر،
وهذا الحد يشمل الأس الصفري أيضًا:
\[
  \fn{Exp}' \ident \lambd[be][\OLId{latin-ordinary}{e} (\fn{Mult}\, \OLId{latin-ordinary}{b}) \num{\OLMathNumeral{1}}].
\]

تنبيه على عبارة الأصل: أطلق الأصل الإنجليزي القول في الحد الأول،
وإنما يصح ذلك مع الأس الموجب. فأما إن كان الأس صفرًا، فحاصل الاختزال
الدالة المحايدة، وليست عدد تشيرش الممثل للواحد؛ فلا يتم بها شرط
التعريف بلامبدا المتقدم. وأما حد الضرب المتكرر فيعطي عدد تشيرش
الممثل للواحد. ونجعل هنا قيمة كل قوة ذات أس صفري واحدًا، ولو كان
أساسها صفرًا، ليعم تعريف دالة القوى جميع أزواج الأعداد الطبيعية.

وأما السابق والطرح على أعداد تشيرش فليسا بهذا اليسر الذي قد يظن؛
وسبيلنا إلى تعريفهما هنا أن نستعمل ترميز الأزواج.

\end{document}
